\documentclass[11pt,a4paper]{article}
\usepackage[utf8]{inputenc}
\usepackage[margin=1in]{geometry}
\usepackage{amsmath, amsfonts, amssymb}
\usepackage{hyperref}

\title{STP-GNN: Technical Foundations and Mathematical Logic}
\author{Yaseen Khalil}
\date{\today}

\begin{document}

\maketitle

This document explains the core components of the STP-GNN framework in simple, rigorous terms. It breaks down the transformation of biological logic into a differentiable machine learning engine.

\section{Algebraic State Space Representation (ASSR)}

\subsection*{What it is and why we need it}
Standard biology models use ``0'' or ``1'' to show if a gene is off or on. However, computers cannot do calculus on single integers. We need to turn these ``bits'' into ``vectors'' so that a gene can be ``80\% on'' or ``20\% off'' during the learning process.

\subsection*{The Math}
For a single gene $x_i$:
\begin{itemize}
    \item On ($1$): $\delta_2^1 = \begin{bmatrix} 1 \\ 0 \end{bmatrix}$
    \item Off ($0$): $\delta_2^2 = \begin{bmatrix} 0 \\ 1 \end{bmatrix}$
\end{itemize}

The global state $X$ for $n$ genes is the multi-index Kronecker product:
\begin{equation}
X = x_1 \otimes x_2 \otimes \dots \otimes x_n
\end{equation}

\subsection*{Explanation}
The symbol $\otimes$ is a Kronecker product. It creates a large vector that represents every possible combination of on/off states for all genes. If you have 10 genes, this vector has $2^{10} = 1024$ positions. Only one position is ``1'' at any time in a pure Boolean system.

\subsection*{Output}
A $2^n$ dimensional vector where a specific index represents the exact ``snapshot'' of the entire cell at that moment.

\section{Temperature-Scaled Softmax}

\subsection*{What it is and why we need it}
In a standard Boolean network, the rules are ``hard'' (e.g., if A is on, B MUST be off). This creates a ``flat'' landscape where the gradient (the slope) is zero. Without a slope, the AI cannot ``climb'' the mountain to find vulnerabilities. We use ``Temperature'' to melt these hard rules into soft curves.

\subsection*{The Math}
\begin{equation}
\tilde{M}_i = \frac{\exp(\tau \cdot \theta_i)}{\sum \exp(\tau \cdot \theta_i)}
\end{equation}
Where $\tau$ (Tau) is our temperature (set to $10.0$ in our optimal build).

\subsection*{Explanation}
As $\tau$ increases, the decision becomes sharper (more like a Boolean gate). By setting it to $10.0$, we keep the gate mostly logical but leave a tiny ``slope'' at the edges. This allows the backpropagation algorithm to ``see'' which way to nudge the parameters to break the system.

\subsection*{Output}
A ``soft'' transition matrix where instead of a $100\%$ chance of a gene turning on, it might be $99.9\%$, leaving $0.1\%$ for the AI to exploit.

\section{Implicit STP (The Khatri-Rao Trick)}

\subsection*{What it is and why we need it}
If we tried to create the full transition matrix for a 10-node network, it would have $1,024 \times 1,024$ entries (1 million). For 20 nodes, it would be 1 trillion. No computer has enough memory for this. We use ``Implicit'' math to calculate the result without ever building the giant matrix.

\subsection*{The Math}
Instead of $X_{t+1} = L \cdot X_t$, we calculate the node-specific results first:
\begin{equation}
y_i = M_i \cdot X_t
\end{equation}
\begin{equation}
X_{t+1} = y_1 \otimes y_2 \otimes \dots \otimes y_n
\end{equation}

\subsection*{Explanation}
We calculate the next state of each gene \textit{individually} first ($y_i$), and then combine them using the Kronecker product. This is mathematically identical to using the giant matrix $L$, but it only requires us to store the small local rules ($M_i$).

\subsection*{Output}
The state of the cell at the next time step ($t+1$) calculated using $1000\times$ less memory.

\section{Adversarial Vulnerability Analysis (PGD)}

\subsection*{What it is and why we need it}
Once the network is differentiable, we act as an ``attacker.'' We want to find the smallest possible change to the biological rules that causes the cell cycle to stop. We use PGD (Projected Gradient Descent) to find these weak points.

\subsection*{The Math}
\begin{equation}
\theta_{new} = \theta_{old} + \epsilon \cdot \text{sign}(\nabla_{\theta} \text{Loss})
\end{equation}

\subsection*{Explanation}
\begin{enumerate}
    \item We define a ``Loss'' (e.g., how well the cell cycle is working).
    \item We calculate the gradient ($\nabla$) to see which rule changes would most effectively \textit{increase} that loss (break the cycle).
    \item We ``nudge'' the rules in that direction by a tiny amount ($\epsilon$).
\end{enumerate}

\subsection*{Output}
A list of ``Topological Vulnerabilities''—the specific genes and conditions that, if perturbed, collapse the system.

\section{Why p53 ``Failed'' vs. Cell Cycle Success}

\subsection*{The Comparison}
\begin{itemize}
    \item \textbf{p53 Model (5 nodes)}: This failed to show ``Holographic'' results because it was too simple. With only 5 genes, the ``Kill Switch'' was too obvious. There weren't enough backup pathways, so any attack worked. It lacked \textit{Topological Complexity}.
    \item \textbf{Cell Cycle Model (10 nodes)}: This was the success. At 10 nodes, the network has enough redundancy that single-gene attacks often fail. This allowed us to prove that you need a \textit{distributed attack} to break a complex system.
\end{itemize}

\subsection*{Why Discrete Logic Fails the Attack}
In a purely discrete simulation, you can only flip a gene from 0 to 1. There is no ``in-between.'' This means you can't use gradients. You would have to try every single combination of flips (brute force), which is impossible for large networks.

\section{DepMap Empirical Validation}

\subsection*{What it is and why we need it}
We need to prove that the ``weak points'' the AI found in our math model actually exist in real human cancer cells. We use the Broad Institute's DepMap data.

\subsection*{The Math (Stratification)}
We take 1,000+ cell lines and group them:
\begin{itemize}
    \item Group A: Genes with working Rb (inhibitor).
    \item Group B: Genes with broken Rb.
\end{itemize}
We then compare their dependency on E2F (the gene our AI attacked).

\subsection*{Explanation}
Our AI said: ``The best way to break the cycle is to exploit the E2F bottleneck.'' 
The DepMap data showed: Cells that already have a broken Rb are \textit{significantly more dependent} on E2F to survive. This means our AI found the exact same biological relationship that exists in real tumors.

\subsection*{Output}
\textbf{Statistical parity}: The AI's predicted ``attack vector'' matches the real-world ``survival dependency'' of cancer cells.

\end{document}
